\documentclass[10pt,a4paper]{article}

% =========================================================
% JOURNAL OF AGRONOMY AND CROP SCIENCE
% Manuscript template inspired by Wiley layout
%
% Required files in Overleaf:
% Journal.png
% original.png
% wiley.png
% ref.bib
% =========================================================

% =========================================================
% ENCODING / LANGUAGE
% =========================================================
\usepackage[T1]{fontenc}
\usepackage[utf8]{inputenc}
\usepackage[english]{babel}

% =========================================================
% FONT / TYPOGRAPHY
% =========================================================
\usepackage{newtxtext,newtxmath}
\usepackage{microtype}
\usepackage{xurl}
\usepackage{graphicx}%
\usepackage{multirow}%
\let\Bbbk\relax
\usepackage{amsmath,amssymb,amsfonts}%
\let\openbox\relax
\usepackage{amsthm}%
\usepackage{mathrsfs}%
\usepackage[title]{appendix}%
\usepackage{xcolor}%
\usepackage{textcomp}%
\usepackage{manyfoot}%
\usepackage{booktabs}%
\usepackage{algorithm}%
\usepackage{algorithmicx}%
\usepackage{algpseudocode}%
\usepackage{listings}%

% For orcid links generation
\RequirePackage{orcidlink}


% =========================================================
% GRAPHICS / COLORS
% =========================================================
\usepackage{graphicx}
\usepackage{xcolor}
% =========================================================
% PAGE GEOMETRY
% =========================================================
\usepackage[
    a4paper,
    top=1.35cm,
    bottom=1.45cm,
    left=1.60cm,
    right=1.60cm,
    headheight=14pt,
    headsep=0.25cm,
    footskip=0.60cm
]{geometry}

% =========================================================
% TEXT / LAYOUT
% =========================================================
\usepackage{ragged2e}
\usepackage{setspace}
\usepackage{titlesec}
\usepackage{enumitem}
\usepackage{calc}

% =========================================================
% TABLES / FIGURES
% =========================================================
\usepackage{booktabs}
\usepackage{multirow}
\usepackage{tabularx}
\usepackage{array}
\usepackage{caption}
\usepackage{dblfloatfix}

% =========================================================
% MATH
% =========================================================
\usepackage{amsmath}

\usepackage{amsfonts}

% =========================================================
% HEADER / FOOTER
% =========================================================
\usepackage{fancyhdr}
\usepackage{lastpage}

% =========================================================
% LINKS
% =========================================================
\usepackage{hyperref}

% =========================================================
% BIBLIOGRAPHY
% =========================================================
\usepackage[
  authordate,
  backend=biber,
  natbib=true,
  doi=true,
  url=false,
  isbn=false,
  giveninits=true,
  uniquename=false,
  maxcitenames=2,
  maxbibnames=99
]{biblatex-chicago}

\defbibenvironment{bibliography}
  {\list{}
    {\setlength{\leftmargin}{0pt}%
     \setlength{\itemindent}{0pt}%
     \setlength{\itemsep}{\bibitemsep}%
     \setlength{\parsep}{\bibparsep}}}
  {\endlist}
  {\item}

\DeclareFieldFormat{doi}{%
  \href{https://doi.org/#1}{https://doi.org/#1}%
}

\addbibresource{ref.bib}

% =========================================================
% GRAPHIC PATHS
% =========================================================
\graphicspath{{./}{figures/}{images/}}

% =========================================================
% GENERAL CONFIGURATION
% =========================================================
\setlength{\parindent}{0pt}
\setlength{\parskip}{0.28em}

\setlength{\columnsep}{0.65cm}
\setlength{\columnseprule}{0pt}

\setlength{\emergencystretch}{2em}

\linespread{0.98}

\raggedbottom

% =========================================================
% COLORS
% =========================================================
\definecolor{JACGray}{RGB}{230,230,230}
\definecolor{JACDark}{RGB}{36,32,33}
\definecolor{JACBlue}{RGB}{0,92,165}
\definecolor{JACRule}{RGB}{75,75,75}

% =========================================================
% HYPERREF
% =========================================================
\hypersetup{
    colorlinks=true,
    linkcolor=JACBlue,
    citecolor=JACBlue,
    urlcolor=JACBlue,
    pdfauthor={Author Name},
    pdftitle={Manuscript Title}
}

% =========================================================
% SECTION TITLES
% =========================================================
\titleformat{\section}
    {\large\bfseries\color{JACDark}}
    {\thesection\,|}
    {0.40em}
    {}

\titleformat{\subsection}
    {\normalsize\bfseries\color{JACDark}}
    {\thesubsection\,|}
    {0.40em}
    {}

\titleformat{\subsubsection}
    {\normalsize\bfseries\color{JACDark}}
    {\thesubsubsection\,|}
    {0.40em}
    {}

\titlespacing*{\section}
    {0pt}
    {0.85em}
    {0.25em}

\titlespacing*{\subsection}
    {0pt}
    {0.60em}
    {0.15em}

\titlespacing*{\subsubsection}
    {0pt}
    {0.50em}
    {0.10em}

% =========================================================
% CAPTIONS
% =========================================================
\captionsetup{
    font=small,
    labelfont=bf,
    labelsep=quad,
    justification=justified,
    singlelinecheck=false
}

% =========================================================
% LISTS
% =========================================================
\setlist[itemize]{
    leftmargin=1.2em,
    itemsep=0.12em,
    topsep=0.15em,
    parsep=0pt
}

\setlist[enumerate]{
    leftmargin=1.4em,
    itemsep=0.12em,
    topsep=0.15em,
    parsep=0pt
}

% =========================================================
% HEADER / FOOTER — PAGE 2 ONWARD
% =========================================================
\pagestyle{fancy}

\fancyhf{}

\renewcommand{\headrulewidth}{0pt}
\renewcommand{\footrulewidth}{0pt}

\fancyhead[L]{%
    \small\itshape
    Journal of Agronomy and Crop Science
}

\fancyhead[R]{%
    \small
    \thepage\ of \pageref{LastPage}
}

\fancyfoot[L]{%
    \small\itshape
    Journal of Agronomy and Crop Science
}

\fancyfoot[R]{%
    \small
    \thepage\ of \pageref{LastPage}
}

% =========================================================
% ARTICLE TITLE
% =========================================================
% USER INPUT: replace the placeholder text below.
% Journal guidance: use a brief, informative title containing the major keywords.
% Avoid abbreviations in the title.
\newcommand{\ArticleTitle}{%
Insert manuscript title here
}

% =========================================================
% AUTHORS
% =========================================================
% Add or remove authors as needed.
% Enter the full names of all authors.
% Use superscript numbers to link authors to affiliations.
% Use * for the corresponding author.
% The journal requires ORCID identifiers.
% Example ORCID format: \orcid{0000-0000-0000-0000}

\newcommand{\orcid}[1]{%
  \hspace{1pt}\orcidlink{#1}\hspace{1pt}%
}

\newcommand{\ArticleAuthors}{%
First Author\orcid{0000-0000-0000-0000}\textsuperscript{1}
\;|\;
Second Author\orcid{0000-0000-0000-0000}\textsuperscript{2}
\;|\;
Corresponding Author\orcid{0000-0000-0000-0000}\textsuperscript{1,*}
}


% =========================================================
% AFFILIATIONS
% =========================================================
% Add, remove, or renumber affiliations as needed.
% Include the institutional affiliation where the work was conducted.
% If an author's present address differs, provide it as a footnote or note as appropriate.
\newcommand{\AffiliationsBlock}{%

\textsuperscript{1}
Department or Laboratory, Institution, City, State/Province, Postal Code, Country.

\vspace{0.05cm}

\textsuperscript{2}
Department or Laboratory, Institution, City, State/Province, Postal Code, Country.

}


% =========================================================
% CORRESPONDENCE
% =========================================================
\newcommand{\CorrespondenceBlock}{%
Correspondence:
Corresponding Author
(\href{mailto:author@example.com}{author@example.com})
}


% =========================================================
% EDITORIAL DATES
% =========================================================
\newcommand{\DatesBlock}{%
Received: -- \;|\;
Revised: -- \;|\;
Accepted: --
}

% =========================================================
% KEYWORDS
% =========================================================
% Journal guidance: provide up to five keywords.
% Include at least one relevant stress topic when applicable to the manuscript.
% Avoid unnecessary duplication of words already used in the title.
\newcommand{\KeywordsBlock}{%
keyword 1 \;|\;
keyword 2 \;|\;
keyword 3 \;|\;
keyword 4 \;|\;
keyword 5
}


% =========================================================
% ABBREVIATIONS
% =========================================================
% If no abbreviations are used, remove this content from the first-page footer.
\newcommand{\AbbreviationsBlock}{%
ABC, abbreviation definition;
XYZ, abbreviation definition.
}


% =========================================================
% JOURNAL INFORMATION
% =========================================================
\newcommand{\RunningYear}{20XX}

\newcommand{\VolumeInfo}{%
Volume:ArticleNumber
}

\newcommand{\DOIInfo}{%
https://doi.org/xx.xxxx/xxxxx
}

% =========================================================
% ABSTRACT
% =========================================================
% Journal guidance: the abstract should be unstructured and should not contain abbreviations.
% Briefly state the context or problem, objective, main methods, principal results, and conclusion.
% Keep the text concise and focused on the study's main contribution.
\newcommand{\AbstractText}{%
\noindent
Insert the abstract here.
}


% =========================================================
% ABSTRACT BOX
% =========================================================
\newcommand{\AbstractBox}[1]{%

\vspace{0.30cm}

\noindent
\colorbox{JACGray}{%
    \parbox{\dimexpr\textwidth-2\fboxsep\relax}{%

        \vspace{0.20cm}

        {\large\bfseries ABSTRACT}

        \vspace{0.10cm}

        \justifying

        #1

        \vspace{0.14cm}

    }%
}

}

% =========================================================
% DOCUMENT
% =========================================================
\begin{document}

% =========================================================
% FIRST PAGE
% =========================================================
\thispagestyle{empty}

% =========================================================
% JOURNAL HEADER
% =========================================================
\noindent
\begin{minipage}[t]{0.45\textwidth}

\vspace{0pt}

\includegraphics[
    width=0.90\linewidth,
    keepaspectratio
]{Journal.png}

\end{minipage}
%
\hfill
%
\begin{minipage}[t]{0.27\textwidth}

\vspace{-0.10cm}

\raggedleft

\includegraphics[
    width=0.95\linewidth,
    keepaspectratio
]{wiley.png}

\end{minipage}

\vspace{0.52cm}

% =========================================================
% ORIGINAL ARTICLE / OPEN ACCESS
% =========================================================
\noindent

\includegraphics[
    width=0.40\textwidth,
    keepaspectratio
]{original.png}

\vspace{0.48cm}

% =========================================================
% TITLE
% =========================================================
{\fontsize{21}{24}\selectfont
\bfseries
\color{JACDark}
\ArticleTitle
\par}

\vspace{0.28cm}

% =========================================================
% AUTHORS
% =========================================================
{\small
\ArticleAuthors
\par}

\vspace{0.18cm}

% =========================================================
% AFFILIATIONS
% =========================================================
{\footnotesize

\setlength{\parindent}{0pt}

\raggedright

\AffiliationsBlock

\par
}

\vspace{0.18cm}

% =========================================================
% CORRESPONDENCE
% =========================================================
{\footnotesize
\CorrespondenceBlock
\par
}

\vspace{0.08cm}

% =========================================================
% DATES
% =========================================================
{\footnotesize

\textbf{Received:} -- 
\;|\;
\textbf{Revised:} --
\;|\;
\textbf{Accepted:} --

\par
}

\vspace{0.08cm}

% =========================================================
% KEYWORDS
% =========================================================
{\footnotesize

\textbf{Keywords:}
\KeywordsBlock

\par
}

% =========================================================
% ABSTRACT
% =========================================================
\AbstractBox{\AbstractText}

\vspace*{\fill}

% =========================================================
% FIRST-PAGE FOOTER
% =========================================================
{\scriptsize

\hrule

\vspace{0.10cm}

\textbf{Abbreviations:}
\AbbreviationsBlock

\vspace{0.10cm}

\hrule

\vspace{0.10cm}

This is an open access article under the terms of the
Creative Commons Attribution License, which permits use,
distribution and reproduction in any medium,
provided the original work is properly cited.

\vspace{0.08cm}

© \RunningYear\ The Author(s).
\textit{Journal of Agronomy and Crop Science}
published by Wiley-VCH GmbH.

\vspace{0.10cm}

\hrule

\vspace{0.10cm}

\textit{Journal of Agronomy and Crop Science},
\RunningYear;
\VolumeInfo

\href{\DOIInfo}{\DOIInfo}

\hfill
1 of \pageref{LastPage}

}

% =========================================================
% PAGE 2 ONWARD — TWO COLUMNS
% =========================================================
\clearpage

\twocolumn

\pagestyle{fancy}

% =========================================================
% KEY POINTS
% =========================================================
% Journal guidance: provide 3 to 5 brief Key Points in bullet form.
% Each point should summarize the initial situation/problem, approach, main result,
% significance, or implication of the research.
% Keep each point brief (approximately 50--100 characters) and avoid abbreviations.
% Replace the example items below with the main take-home messages of the manuscript.
\noindent
\fbox{%
\begin{minipage}{0.94\columnwidth}

\vspace{0.14cm}

{\normalsize\bfseries Key Points}

\vspace{0.06cm}

\begin{itemize}
    \setlength{\itemsep}{0.05cm}
    \setlength{\topsep}{0.02cm}
    \setlength{\parsep}{0pt}
    \setlength{\partopsep}{0pt}

    \item Insert the first key point here.
    \item Insert the second key point here.
    \item Insert the third key point here.
    \item Insert the fourth key point here.
    \item Insert the fifth key point here.

\end{itemize}

\vspace{0.04cm}

\end{minipage}%
}

\vspace{0.22cm}

% =========================================================
% 1 INTRODUCTION
% =========================================================
\section{Introduction}
\label{sec:introduction}

% Suggested content:
% - Introduce the broader scientific context and the relevance of the topic.
% - Define the specific problem or knowledge gap addressed by the study.
% - Summarize the most relevant literature needed to justify the research.
% - Explain why the study is needed and what is novel about the approach.
% - End the section with a clear objective and, when appropriate, the hypothesis.
% Use \citet{key} for narrative citations and \citep{key} for parenthetical citations.


% =========================================================
% 2 MATERIALS AND METHODS
% =========================================================
\section{Materials and Methods}
\label{sec:methods}

% Adapt the subsection titles below to the manuscript.
% Provide enough methodological detail for the study to be understood and reproduced.
% Describe materials, experimental or study design, data collection, and analytical procedures.
% Include ethical approvals, permits, or reporting standards here when applicable.

\subsection{Study Material or Experimental Material}
\label{subsec:material}

% Describe the biological material, germplasm, study population, experimental units,
% locations, growing seasons, or other materials relevant to the study.

\subsection{Experimental Design or Study Design}
\label{subsec:design}

% Describe the experimental layout or study design, treatments, replication, randomization,
% management practices, and other design features needed to interpret the study.

\subsection{Data Collection}
\label{subsec:data}

% Describe the traits, variables, measurements, sampling procedures, instruments,
% environmental data, laboratory procedures, or other data sources used in the study.

\subsection{Statistical Analysis}
\label{subsec:statistics}

% Describe all statistical models, software, packages, significance criteria,
% validation procedures, and any data transformations or quality-control steps.
% Provide software versions and citations when appropriate.


% =========================================================
% FIGURE — TWO COLUMNS
% =========================================================
% Example:
% \begin{figure*}[t]
% \centering
% \includegraphics[width=0.96\textwidth]{figure1.png}
% \caption{Insert figure caption here.}
% \label{fig:figure1}
% \end{figure*}

% =========================================================
% FIGURE — ONE COLUMN
% =========================================================
% Example:
% \begin{figure}[t]
% \centering
% \includegraphics[width=\linewidth]{figure2.png}
% \caption{Insert figure caption here.}
% \label{fig:figure2}
% \end{figure}

% =========================================================
% TABLE — TWO COLUMNS
% =========================================================
% Example:
% \begin{table*}[t]
% \centering
% \caption{Insert table caption here.}
% \label{tab:table1}
% \begin{tabular}{llll}
% \toprule
% Column 1 & Column 2 & Column 3 & Column 4 \\
% \midrule
% -- & -- & -- & -- \\
% \bottomrule
% \end{tabular}
% \end{table*}


% =========================================================
% 3 RESULTS
% =========================================================
\section{Results}
\label{sec:results}

% Present the main findings in a clear and logical order.
% Report results objectively and support them with the relevant tables and figures.
% Avoid extensive interpretation here; reserve broader interpretation for the Discussion.
% Add or remove subsections as needed.

\subsection{First Results Subsection}
% Add text here.

\subsection{Second Results Subsection}
% Add text here.


% =========================================================
% 4 DISCUSSION
% =========================================================
\section{Discussion}
\label{sec:discussion}

% Suggested content:
% - Interpret the main findings in relation to the study objective or hypothesis.
% - Compare the results with relevant previous studies.
% - Explain biological, agronomic, methodological, or practical implications.
% - Discuss important limitations without overstating the conclusions.
% - Identify relevant directions for future research when appropriate.


% =========================================================
% 5 CONCLUSION
% =========================================================
\section{Conclusion}
\label{sec:conclusion}

% Summarize the principal conclusions supported by the results.
% State the main contribution and practical or scientific implications of the study.
% Avoid introducing new results, citations, or unsupported generalizations.

% =========================================================
% AUTHOR CONTRIBUTIONS
% =========================================================
\section*{Author Contributions}

% The journal requires author contributions to follow the CRediT taxonomy.
% List each author's initials or name followed by the relevant contribution roles.
% Example: A.A.: Conceptualization, Methodology, Writing -- original draft.
% B.B.: Data curation, Formal analysis, Writing -- review and editing.


% =========================================================
% ACKNOWLEDGEMENTS
% =========================================================
\section*{Acknowledgements}

% Acknowledge individuals or organizations that contributed to the work but do not meet authorship criteria.
% The journal also asks authors to list funding sources accurately; funding may be stated here
% and/or in the dedicated Funding section below, depending on the final submission format.


% =========================================================
% FUNDING
% =========================================================
\section*{Funding}

% List all funding sources and grant numbers accurately.
% Use the official funder name whenever possible.
% If no specific funding was received, include an appropriate no-funding statement.


% =========================================================
% CONFLICT OF INTEREST
% =========================================================
\section*{Conflict of Interest}

% Disclose any financial or non-financial interests that could be perceived as influencing the work.
% If there are no conflicts of interest, include an explicit declaration.


% =========================================================
% DATA AVAILABILITY
% =========================================================
\section*{Data Availability Statement}

% The journal expects a Data Availability Statement.
% State where the data supporting the findings can be accessed, including repository name, DOI,
% accession number, or URL when available. If data are restricted, explain the reason and access conditions.


% =========================================================
% SUPPLEMENTARY INFORMATION
% =========================================================
%\section*{Supplementary Information}

% Describe any supplementary files provided with the manuscript.
% Supporting information should be supplied as separate files at submission.




% =========================================================
% REFERENCES
% =========================================================
\section*{References}
% Keep the reference list consistent. The journal uses APA style, but Free Format submission
% allows references to be submitted in another consistent style at initial submission.
\printbibliography[heading=none]


% =========================================================
% END DOCUMENT
% =========================================================
\end{document}